\chapter{Trabalhos relacionados}
\label{cap:relacionados}

Este capítulo aproxima a proposta de quatro linhas de pesquisa: índices construídos com notícias, modelos de sentimento financeiro em português, análise orientada a entidades e estudos que relacionam texto e mercado brasileiro. O objetivo não é transportar resultados entre bases diferentes, mas identificar quais decisões metodológicas foram adotadas na literatura e em que aspectos o ITI se aproxima ou se distingue delas.

\section{Indicadores construídos a partir de notícias}

\citeonline{tetlock2007media} utilizou linguagem da imprensa para investigar sua relação com o mercado acionário norte-americano. Esse estudo demonstra que o conteúdo jornalístico pode ser representado quantitativamente, desde que sejam identificados o tipo de texto, o fenômeno medido e o desfecho financeiro utilizado. Ele não oferece, entretanto, um procedimento pronto para classificar notícias brasileiras em relação a uma empresa específica.

O estudo de \citeonline{baker2016epu} construiu um índice de incerteza de política econômica que utiliza, entre outros componentes, a cobertura jornalística. Sua contribuição para esta dissertação está em mostrar que a seleção de textos precisa corresponder ao construto de interesse. Uma notícia pode ser relevante para medir incerteza sem apresentar orientação positiva ou negativa identificável para a Sabesp.

\citeonline{shapiro2022fed} desenvolveram medidas de sentimento a partir de notícias econômicas e financeiras e avaliaram modelos de classificação com apoio de anotações humanas. O \emph{Daily News Sentiment Index} do Federal Reserve Bank of San Francisco utiliza agregação diária e ponderação de notícias anteriores com pesos decrescentes \cite{frbsfNewsSentiment,buckman2020dnsi}. Esse antecedente é especialmente próximo da ideia de preservar informações no tempo, mas seu alvo econômico, seleção de fontes, ajuste de composição e regra de agregação não são idênticos aos do ITI.

\citeonline{hiew2019bsi} também estudaram a construção de índice de sentimento financeiro a partir de modelos de linguagem, relacionando a medida a dados de mercado. O trabalho sustenta a investigação de representações agregadas que não se limitam à análise individual de cada texto. Nesta pesquisa, um contraste por escore discreto, denominado Hard-BSI na implementação, serve como referência de agregação por classe prevista; ele não reproduz integralmente o método daquele estudo.

\section{Modelos de sentimento financeiro em português}

No domínio financeiro, o significado de palavras pode depender do contexto econômico. \citeonline{loughran2011liability} mostraram limitações de abordagens lexicais gerais aplicadas a textos financeiros; \citeonline{araci2020finbert} investigou o emprego de modelos BERT nesse domínio. Esses trabalhos ajudam a justificar a escolha de um classificador ajustado ao vocabulário financeiro em lugar de uma abordagem exclusivamente baseada em palavras isoladas.

No Brasil, \citeonline{santos2023finbert} apresentaram o FinBERT-PT-BR e examinaram sua utilização em textos financeiros em português. Os resultados do estudo original, incluindo aproximadamente 0,76 de acurácia e 0,73 de F1 nas condições relatadas pelos autores, fornecem um antecedente para a escolha do modelo. Não devem, contudo, ser interpretados como desempenho esperado em notícias extensas de saneamento ou como garantia de sentimento orientado à Sabesp.

O mesmo estudo explora a agregação de classificações em uma medida quantitativa baseada na contagem de textos positivos e negativos. O ITI difere desse procedimento ao utilizar a diferença entre as probabilidades de cada notícia, calcular médias diárias e aplicar memória exponencial. Essas escolhas são formuladas e avaliadas no presente trabalho; não são apresentadas como resultados diretamente demonstrados pelo artigo do FinBERT.

\section{Sentimento orientado a entidades}

A tarefa SemEval-2017 Task 5, apresentada por \citeonline{cortis2017semeval5}, investigou a análise fina de sentimento em textos financeiros com referências explícitas a empresas e instrumentos. Em notícias sobre organizações distintas, atribuir um único rótulo ao documento pode ser insuficiente quando os efeitos econômicos descritos não são os mesmos para cada entidade.

\citeonline{yuan2020target} estudaram a anotação de notícias financeiras chinesas com alvos definidos. \citeonline{sinha2022sentfin}, por sua vez, apresentaram o corpus SEntFiN 1.0, com identificação da orientação de manchetes em relação às entidades mencionadas. \citeonline{tang2023finentity} também investigaram a classificação de sentimento em nível de entidade. Em conjunto, esses trabalhos mostram que identificar o alvo não é um detalhe secundário da tarefa, mas uma decisão necessária para interpretar seus resultados.

A presente pesquisa utiliza essa literatura para fundamentar a pergunta feita aos anotadores e a preparação da entrada textual próxima às menções à Sabesp. Existe, porém, uma distinção importante: selecionar sentenças sobre a empresa não equivale a treinar um modelo para reconhecer explicitamente o par notícia--entidade. A capacidade do FinBERT-PT-BR de atender à tarefa ainda precisa ser medida na entrada efetivamente utilizada pelo ITI.

\section{Investigações no mercado brasileiro}

\citeonline{yoshinaga2012sentimento} analisaram relações entre sentimento de mercado e retornos acionários no Brasil. Seu índice foi construído a partir de variáveis diferentes daquelas utilizadas neste trabalho; por isso, os resultados não são diretamente comparáveis, mas reforçam a importância de interpretar o sinal e a natureza da medida investigada.

\citeonline{duarte2020news} estudaram a utilização de notícias para investigar quedas de preços no mercado brasileiro. O trabalho aproxima processamento textual e desfecho financeiro, mas sua tarefa de previsão de eventos não equivale a uma comparação de correlações marginais entre séries semanais. \citeonline{racef2022sentimento} também examinaram associação entre sentimento de investidores e fases do mercado, lembrando que o período analisado importa para a interpretação de relações financeiras.

A reorganização societária da Sabesp constitui parte do contexto institucional do setor. O estudo de \citeonline{foco2025sabesp}, publicado em 2026, examinou, por estudo de eventos, a reação das ações da Equatorial Energia (EQTL3) a um anúncio ligado à privatização da Sabesp. Seu desfecho é diferente do retorno da SBSP3 utilizado nesta dissertação, e a publicação não valida a construção do ITI. Ela é citada apenas como exemplo de investigação de um acontecimento relacionado ao setor com outro método e outra variável financeira.

\section{Posicionamento da pesquisa}

A literatura oferece fundamentos para quatro decisões: transformar textos em medidas, adotar um classificador financeiro em português, considerar a empresa-alvo e estudar o papel de notícias anteriores por meio de agregação temporal. A contribuição do ITI nesta etapa está em integrar essas decisões num procedimento explícito e submetê-lo a verificações textuais e temporais no estudo de caso da Sabesp. A Tabela~\ref{tab:posicionamento} resume o posicionamento, sem alegar que os procedimentos anteriores são equivalentes.

\begin{table}[htbp]
\centering
\caption{Relação entre a proposta do ITI e linhas de pesquisa anteriores.}
\label{tab:posicionamento}
\begin{tabularx}{\textwidth}{@{}P{0.22\textwidth}L L@{}}
\toprule
Linha & Antecedentes & Delimitação neste trabalho \\
\midrule
Notícias e mercado & \citeonline{tetlock2007media}; \citeonline{baker2016epu} & Corpus em português e alvo empresarial definido \\
Índices de notícias & \citeonline{shapiro2022fed}; \citeonline{buckman2020dnsi} & Média diária e meia-vida em dias civis, sem equivalência metodológica presumida \\
Sentimento financeiro & \citeonline{santos2023finbert} & Probabilidades e avaliação local do classificador \\
Sentimento por entidade & \citeonline{cortis2017semeval5}; \citeonline{sinha2022sentfin}; \citeonline{tang2023finentity} & Protocolo orientado à Sabesp e recorte contextual, ainda sem modelo especializado por alvo \\
Mercado brasileiro & \citeonline{yoshinaga2012sentimento}; \citeonline{duarte2020news} & Associação exploratória com retornos futuros de SBSP3 \\
\bottomrule
\end{tabularx}
\vspace{0.4em}
{\footnotesize Fonte: elaboração própria com base nos trabalhos citados.}
\end{table}
